% 编译：latexmk -xelatex main.tex
% 新报告：复制整个文件夹，只改本文件和 refs.bib。
\documentclass{csee-report}

% ---------------- 首页编号行（可选，各项都不写就不排） ----------------
% \clc{TM 614}          % 中图分类号，字母与数字间一个空格
% \doccode{A}           % 文献标志码
% \subjectcode{470.40}  % 学科分类号
% 文章编号 \articleno{} 由期刊分配，课程报告不写

% ---------------- 题名页 ----------------
\title{风电场自主运行控制方法综述}
\author{张三，李四，王五}
\affil{华中科技大学电气与电子工程学院，湖北省 武汉市 430074}
\entitle{A Review of Autonomous Operation and Control Methods for Wind Farms}
\enauthor{ZHANG San, LI Si, WANG Wu}
\enaffil{School of Electrical and Electronic Engineering, Huazhong University of Science and Technology, Wuhan 430074, Hubei Province, China}

\enabstract{English abstract here, 9 pt Times New Roman with 14 pt line spacing.
It should state the objective, method, results and conclusion of the report.}
\enkeywords{wind farm; autonomous operation; flow control; lifetime control}
\abstract{中文摘要，小五号宋体，行距 14 pt。摘要应说明研究目的、方法、结果和结论，
不写“本文”以外的背景铺垫。}
\keywords{风电场；自主运行；尾流控制；寿命控制}
% \fund{国家自然科学基金项目(00000000)。}   % 有基金项目时取消注释，排在首页左栏底部
% \oddheader{张三等：风电场自主运行控制方法综述}  % 奇数页页眉居中内容，默认为空
% 作者简介（可选，排在参考文献之后）：\biography{照片文件}{姓名}{简介}，照片放 figures/，宽 2.2 cm、高 2.8 cm
% \biography{zhangsan.jpg}{张三}{张三(1990)，男，硕士研究生，主要研究方向为风电场控制；zhangsan@example.com。}

\begin{document}

\maketitle

\section{引言}
正文五号宋体，英文和数字为 Times New Roman，首行缩进 2 字符，单倍行距(每页 42 行)。
括号除列项编号“1）”外一律用半角。
国内首个“无人化”风电场已实现巡检运维由智能设备承担\cite{zhao2025wuren}，
下一步是让风电场在异常工况下自主决策\cite{barker2020resilient,meyers2022flow}。
数值与单位之间用 \verb|\SI|，如 \SI{12345.87908}{\micro\metre}、\SI{70}{MW}。

文中列项格式如下：
\begin{enumerate}
  \item 工作模式 1。
  \item 工作模式 2：列项较长时，换行后顶格排，与机报模板的“1）”列项格式一致。
\end{enumerate}

\section{风电机组寿命控制}
\subsection{基于预测的自适应控制}
文献\cite{kipchirchir2023prognostics}根据剩余寿命预测调整控制策略。公式右侧编号，如式\eqref{eq:power}：
\begin{equation}
  P = \frac{1}{2}\rho \pi R^{2} C_\mathrm{p}(\lambda,\beta) v^{3}
  \label{eq:power}
\end{equation}
式中：$\rho$ 为空气密度；$R$ 为风轮半径；$C_\mathrm{p}$ 为风能利用系数；$v$ 为风速。

多行公式用 \verb|align|，每行编号，\verb|&| 处对齐；不需要编号的行加 \verb|\nonumber|：
\begin{align}
  \lambda &= \frac{\omega_\mathrm{r} R}{v} \label{eq:tsr} \\
  T_\mathrm{m} &= \frac{P}{\omega_\mathrm{r}}
\end{align}

\subsubsection{四级标题}
四级标题为五号宋体，与正文接排在下一行。图题、表题用 \verb|\bicaption{中文}{English}|，
图题在图下，表题在表上，见图\ref{fig:example}和表\ref{tab:example}。

\begin{figure}[htbp]
  \centering
  \includegraphics[width=0.85\columnwidth]{example-image-a}
  \bicaption{示例图片}{Example figure}
  \label{fig:example}
\end{figure}

\begin{table}[htbp]
  \centering
  \bicaption{DPMA 与 DFT 算法的性能比较}{Performance comparison between DPMA and DFT algorithm}
  \label{tab:example}
  \begin{fulltabular}{ccccc}
    \toprule
    算法 & $f_\mathrm{sp}$/Hz & $e_\mathrm{ang}$/(°) & $e_\mathrm{mag}$/‰ & $e_\mathrm{TVE}$/\% \\
    \midrule
    \multirow{2}{*}{DFT}  & 2 & \num{0.0997} & \num{2.1959} & \num{9.9716} \\
                          & 3 & \num{0.1486} & \num{3.6945} & \num{14.8633} \\
    \midrule
    \multirow{2}{*}{DPMA} & 2 & \num{0.0073} & \num{1.2410} & \num{1.2381} \\
                          & 3 & \num{0.0173} & \num{2.7717} & \num{2.7751} \\
    \bottomrule
  \end{fulltabular}
\end{table}

\section{结论}
总结报告的主要结论。

\begin{acknowledgment}
致谢内容(可删去)。
\end{acknowledgment}

\bibliography{refs}

\appendix
\section{}
附录中的图、表、公式编号为 A1、A2……
\begin{equation}
  \lambda = \frac{\omega R}{v}
\end{equation}

\end{document}
